\chapter{Metodologia}
\label{cap:metodologia}

% Chris Eng.ish --- III. Methodology
% Materials used | Design of the Artefact | Use of the Artefact

\section{Materiais utilizados}
\label{sec:materiais}

\begin{description}
  \item[Laboratório] CONCEPTIO --- drones reais; broker MQTT do lab; protocolo
        de telemetria padronizado.
  \item[Modelo SysML~v2] Editor SysON (Docker local) como anfitrião do modelo no
        PoC; marcador \texttt{\_go\_to\_verification}; Success Criteria como
        atributos/itens no modelo; resultado em \texttt{VerificationResult}.
  \item[Software] ReqValLive (FastAPI + UI); schema Vampire
        (\texttt{simreqvalidator}); opcionalmente LLM (Ollama) para enriquecimento
        de narrativa pós-medição.
  \item[Referências de V\&V] SIS-08 Methods; MSFC-HDBK-3173~\cite{msfc3173}.
\end{description}

% CATIA Magic 2026x permanece contexto do lab / evolução TWC; não é runtime live
% da PoC de medição.

\section{Desenho do artefacto}
\label{sec:desenho-artefacto}

% Design of the Artefact --- método + arquitectura ReqValLive.

\subsection{Ciclo operacional proposto}

\begin{enumerate}
  \item Formalizar requisito + Success Criterion no modelo (SysON / SysML~v2).
  \item Importar e aplicar gate ACCEPT/REJECT antes da medição.
  \item Congelar SC aprovado; subscrever MQTT; avaliar live.
  \item Encerrar; gerar laudo HTML; publicar \texttt{VerificationResult} no modelo
        (e marcadores textuais quando aplicável).
\end{enumerate}

\subsection{Gate de Success Criteria}

Antes da run, cada critério é classificado ACCEPT ou REJECT com motivos
explícitos (completude, mensurabilidade, unidade, escopo, coerência Method ×
Success, temporalidade).

\subsection{Avaliação temporal live}

Critérios com sentido temporal (ex.: \texttt{all\_timesteps}) não podem ser
julgados apenas pelo último valor. A evidência preserva violação (latch FAIL),
primeira falha, extremos e contagens.

O vocabulário de observação segue a taxonomia GSE de Jefferson
(sete naturezas): valor, tempo, sequência, taxa, ocorrência, frequência e
consistência de estados~\cite{jefferson2026gse}.

Tipos suportados na PoC: \texttt{threshold} e \texttt{range} (valor);
\texttt{statistical} com agregações \texttt{range}/\texttt{max}/\texttt{min}
(amplitude na janela, sem $\Delta t$);
\texttt{rate} (taxa $\Delta/\Delta t$);
\texttt{dwell} (a condição só falha se persistir $\geq T$\,ms);
métricas agregadas entre entidades (ex.: separação mínima).
Sequência, eventos digitais, frequência e invariantes de estado
permanecem no backlog.

\subsection{Arquitectura da prova de conceito (ReqValLive)}

Aplicação web local: sessão de medição, subscriber MQTT multi-drone, motor de
avaliação, relatório HTML, integração SysON (import SC + publish
\texttt{VerificationResult}).

\subsection{Artefatos gerados}

\begin{itemize}
  \item Relatório HTML (procedimento de V\&V: método, SC aprovado, esperado vs
        observado, evidência).
  \item Item \texttt{VerificationResult} ligado ao requisito no SysON (status,
        reason, extremos, sessão).
  \item Esqueleto \texttt{.sysml} / model.md como ponto de partida MBSE --- não
        simulação Magic.
\end{itemize}

\subsection{Limitações de desenho assumidas}

Não se pretende adoção industrial, plataforma WARA-PS, nem substituir ferramentas
MBSE comerciais. O modelo exportado não é simulação funcional no CATIA Magic
2026x.

\section{Uso do artefacto}
\label{sec:uso-artefacto}

% Use of the Artefact --- fluxo do operador / investigador.

\subsection{Fluxo típico no laboratório}

\begin{enumerate}
  \item Subir SysON + ReqValLive; garantir MQTT do lab acessível.
  \item Autorar/importar requisitos com SC e \texttt{\_go\_to\_verification}.
  \item Importar no app; resolver gate; montar GSE / sessão.
  \item Conectar MQTT; iniciar medição; encerrar.
  \item Abrir laudo HTML; publicar resultado no SysON; refrescar explorador.
\end{enumerate}

\subsection{Papéis}

\begin{itemize}
  \item \textbf{Investigador / mestrando:} define SC, conduz ensaio, interpreta
        laudo.
  \item \textbf{Modelo (SysON):} autoridade do requisito e do resultado escrito
        de volta.
  \item \textbf{ReqValLive:} motor de evidência live (não substitui o modelo).
\end{itemize}
